\section{Introducción}
La incertidumbre del pasajero—no saber por dónde va el bus o cuánto falta—se reduce al mostrar posiciones en vivo, algo reportado en el transporte urbano por \cite{cityruta2018new} y \cite{stopbus2016new}. UrbanTracker busca eso y, además, facilitar la operación diaria con un panel para rutas, vehículos y conductores.

Para evitar costos de entrada altos, aprovechamos el teléfono del conductor como emisor de GPS cuando el bus no trae equipo dedicado; este enfoque ya ha sido útil en rastreo y seguridad vehicular con mensajería ligera MQTT \cite{stopbus2016new}; \cite{sismo2016}. En paralelo, una API REST clara y documentada evita malentendidos y acelera integraciones \cite{easyrest2022}; \cite{ingenieria2022}.

En móviles escogimos React Native por experiencia del equipo, aunque Flutter es una alternativa sólida de alto rendimiento en UI; conviene evaluarla cuando el proyecto crece o cambian los requisitos \cite{flutter2022new}.